% !PRINTED 538
\weberpage{538}

If the $y$ are then determined in this way, we set
\weq{10}{\omega = \omega_1 y_1 + \omega_2 y_2 + \cdots + \omega_n y_n,}
and derive from (5) the equation
\weq{11}{\lambda\,\omega = \sum^{s} \alpha_s u_s.}

Now since the $\alpha_s$ are divisible by $\lambda$ and the $u_s$ by $p$,
it follows that $\omega$ is divisible by $p$, and hence, by
\S.\,145, 2., that $y_1, y_2, \ldots, y_n$ must be divisible by $p$, which
is contrary to our assumption.

This proves that $T$ is a primitive function, \gesperrt{and
that therefore the absolute value of the determinant $A$
is equal to the absolute norm of $\lambda$, and consequently also
to that of $\mu$} (\S.\,138).

If we apply the result to the special basis (3), \S.\,146,
we obtain the formula:
\weq{12}{N_a(\mu) = a_{1,1}\,a_{2,2}\ldots a_{n,n}.}

To this we append the following remarks: If
in a basis form of a functional $\mu$
\[
\lambda = \alpha_1 t_1 + \alpha_2 t_2 + \cdots + \alpha_n t_n
\]
one applies to the variables $t_1, t_2, \ldots, t_n$ a linear
substitution with integral coefficients and determinant $\pm 1$, a
new basis form of $\mu$ arises. For to all integral rational values
of the variables $t_1, t_2, \ldots, t_n$ there correspond integral rational
values of the new variables, and conversely.

But the application of a linear substitution to the variables $t$
is equivalent to the application of the transposed
substitution to the coefficients $\alpha_1, \alpha_2, \ldots, \alpha_n$ (\S.\,37), that is, to
the transition to a new basis of $\mu$:
\weq{13}{(\beta_1, \beta_2, \ldots, \beta_n) = (B)\,(\alpha_1, \alpha_2, \ldots, \alpha_n),}
which then, through composition with (1), yields:
\weq{14}{(\beta_1, \beta_2, \ldots, \beta_n) = (B)\,(A)\,(\omega_1, \omega_2, \ldots, \omega_n).}

If the discriminant (\S.\,144) of the numbers $\beta$ determined
by a substitution (13) agrees with the discriminant of the $\alpha$,
then the substitution determinant must be $B = \pm 1$;
and we thus arrive at the theorems:
\begin{enumerate}
\item[1.] \gesperrt{The discriminant of a basis of $\mu$ is equal
to the square of the absolute norm of $\mu$,
multiplied by the fundamental number of the field;}
\end{enumerate}
%CONT>>
